\section{Pendahuluan}

Sistem tanya jawab berbasis \textit{Retrieval-Augmented Generation} (RAG) memerlukan sumber eksternal agar jawaban tidak hanya bergantung pada pengetahuan parametrik model bahasa \cite{lewisRetrievalAugmentedGenerationKnowledgeIntensive2021,huangSurveyRetrievalAugmentedText2024}. Kualitas sumber eksternal tersebut ditentukan jauh sebelum model menghasilkan jawaban: dokumen harus diubah menjadi unit yang dapat dicari, indeks harus memelihara identitas sumber, dan hubungan antardokumen perlu direpresentasikan ketika jawaban memerlukan konteks lintas dokumen. Tanpa ketiga hal tersebut, sistem dapat mengembalikan potongan yang benar secara leksikal tetapi salah konteks.

Kebutuhan tersebut menjadi lebih sulit pada proyek Capstone yang menggabungkan empat kontribusi teknis ke dalam satu platform. Proyek ini membangun OMNI-RAG, yaitu platform RAG multidomain untuk chatbot asisten hukum peraturan perundang-undangan Indonesia. Korpus hukum bersifat hierarkis: Bab, Bagian, Pasal, Ayat, Huruf, dan Angka mempunyai hubungan induk-anak yang memengaruhi makna ketentuan \cite{uuPembentukanPeraturan2011}. Di sisi lain, platform juga harus dapat memproses dokumen non-hukum seperti \textit{User Manual} yang memiliki struktur berbeda. Dengan demikian, satu strategi pemotongan atau pencarian tidak dapat diasumsikan cocok untuk semua sumber.

Paper ini menyajikan kontribusi individual pada pembentukan basis pengetahuan OMNI-RAG. Fokusnya adalah tiga komponen pada jalur offline, yaitu \chunker, \indexer, dan pembangun \kg{}, serta mekanisme \plugin{} dan \profile{} yang mengikat ketiganya menjadi pipeline yang dapat dipertukarkan. Jalur offline menerima keluaran parser berupa dokumen kanonik, lalu menghasilkan unit pencarian, indeks, dan graf yang digunakan oleh komponen retrieval pada jalur online. Komponen akuisisi, antarmuka percakapan, dan lapisan aplikasi lain dipaparkan hanya sebagai konteks integrasi.

Pertanyaan penelitian paper ini adalah sebagai berikut.
\begin{enumerate}
  \item Bagaimana kontrak artefak kanonik dan arsitektur \plugin{} dapat memisahkan strategi pemotongan, pengindeksan, dan pembentukan graf dari orkestrator inti?
  \item Bagaimana strategi \textit{structure-aware} dibandingkan dengan strategi generik pada dokumen manual dan peraturan hukum?
  \item Bagaimana sparse, dense, dan hybrid \indexer{} memengaruhi kelengkapan dan urutan konteks yang ditemukan?
  \item Sejauh mana graf hukum dapat memperluas lingkungan dokumen yang ditemukan tanpa mengubah metrik retrieval langsung?
\end{enumerate}

Kontribusi paper ini meliputi: (1) rancangan kontrak \canon{CanonicalDocument}, \canon{CanonicalChunks}, \canon{CanonicalIndex}, dan \canon{CanonicalGraph} yang dipakai lintas domain; (2) implementasi dan evaluasi \plugin{} untuk \chunker, \indexer, dan \kg{} dengan pemilihan melalui \profile{}; (3) prosedur profil, pipeline run, attempt, dan provenance untuk membandingkan eksperimen secara terisolasi; serta (4) analisis hasil pada dua korpus yang menunjukkan bahwa struktur dokumen dan tujuan metrik menentukan konfigurasi terbaik.

Ruang lingkup paper dibatasi pada konstruksi dan evaluasi basis pengetahuan. Pembahasan tidak mengklaim efektivitas seluruh chatbot atau kualitas sistem di luar konfigurasi eksperimen yang dilaporkan. Bagian selanjutnya menjelaskan konteks Capstone dan rancangan arsitektur, diikuti metode eksperimen, hasil, keterbatasan, dan kesimpulan.

